\chapter{附录：时间线、清单与关键文件}
\label{app:tools}
\ist{时间线}

\section{全国统一时间轴（2027 届，2026 年秋季）}

下面这条时间轴来自教育部推免服务系统的官方说明与省级教育考试机构的公告，属于\textbf{全国统一节点}。推免系统的时间安排逐年可能微调，\textbf{每年务必以研招网当年公告为准}。\nref{S0464}, \nref{S0389}

\begin{timeline}
\tl{2026-09-18}{推免服务系统开放注册，推免生填报个人资料、查看本人推免资格}
\tl{2026-09-21 09:00}{开始填报志愿（至 10-18 截止填报）}
\tl{2026-09-22 08:00}{开始确认招生单位发送的复试通知}
\tl{2026-09-24 14:00}{开始确认招生单位发送的待录取通知}
\tl{2026-10-13 至 10-20}{2023 级拟转段强基生查询并确认本人转段信息}
\tl{2026-10-18}{截止填报志愿}
\tl{2026-10-20}{截止确认复试及待录取通知}
\end{timeline}


\subsection*{系统层面的三条硬规则}

\begin{itemize}
\item 推免生（有特殊政策要求的专项计划推免生除外）\textbf{可同时填报三个平行志愿}，不分主次。
\item 每个志愿\textbf{提交后 48 小时内不允许修改}。若招生单位在 48 小时内通过系统明确拒绝，推免生可立即填报其他志愿；48 小时后仍未接到复试通知、复试未通过，或推免生拒绝待录取通知的，该志愿自动解锁。
\item \textbf{只能接受一个待录取通知}。未在系统中确认接受待录取通知的，录取不生效。
\end{itemize}

\nref{S0236}, \nref{S0444}, \nref{S0182}

\subsection*{全国统一的制度性约束}

\begin{itemize}
\item 各招生单位\textbf{接收推免硕士生数量不得超过本单位全日制硕士研究生招生计划的 50\%}，各全日制专业必须留出一定比例用于考试招生。
\item 拟推免名单、拟接收名单等重要事项须统一公示，\textbf{公示时间不少于 10 个工作日}；公示期间内容不得修改，如有变动须另行公示 10 个工作日。
\item 推免名额比例（2006 年试行办法）：设研究生院高校约为应届本科毕业生数的 15\%，未设研究生院的“211 工程”高校约 5\%，其他高校约 2\%。
\item 截至规定时间被确定的推免生若无招生单位接收或本人放弃，\textbf{其推免资格作废，名额不得转让}。
\end{itemize}

\nref{S0318}, \nref{S0047}, \nref{S0389}

\begin{keybox}
\textbf{注意比例文件的年代。} 上述 15\%/5\%/2\% 出自 2006 年《推免工作管理办法（试行）》，该办法的框架自 2026 年起已被新的《全国高等学校推荐优秀应届本科毕业生免试攻读研究生工作管理规定》取代。\textbf{引用具体比例前，请以本校当年推免办法为准}——各校实际名额由教育部逐年下达。
\end{keybox}

\section{关键节点自查清单}

\subsection*{阶段一：资格确认（大三下 \textasciitilde{} 大四上）}

\begin{itemize}
\item[$\square$] 查清本校推免办法中\textbf{成绩排名的计算口径}（加权平均分／学分绩点／算术平均；含不含选修；重修如何计）
\item[$\square$] 确认\textbf{计入排名的学期范围}（前五学期／前六学期），这直接决定夏令营能否投递
\item[$\square$] 核对是否有\textbf{挂科、处分、体测}等一票否决或减分项
\item[$\square$] 核对加分项认定目录与\textbf{加分上限}（论文、竞赛、专利、社会工作等）
\item[$\square$] 确认名额分配方式（按专业人数比例／学院切块／学校统筹）
\item[$\square$] 关注\textbf{公示期}与异议申诉渠道，公示期内发现问题及时提出
\end{itemize}

\subsection*{阶段二：申请与考核（6 \textasciitilde{} 8 月）}

\begin{itemize}
\item[$\square$] 逐校核对目标学院的\textbf{申请条件}（排名门槛、英语门槛、是否限本科出身、是否限专业）
\item[$\square$] 按该院当年简章准备材料清单，注意\textbf{推荐信份数与推荐人职称要求}、个人陈述字数上限
\item[$\square$] 确认报名系统（学校自有系统 \emph{与} 教育部推免系统是否需要\textbf{两段都走}）
\item[$\square$] 确认考核形式（面试／笔试／机试；线上／线下）与分值构成
\item[$\square$] 弄清\textbf{“优秀营员”在本院的实际效力}（可直接拟录取／仅优先／纯荣誉），务必看简章原文
\end{itemize}

\subsection*{阶段三：九推与确认（9 \textasciitilde{} 10 月）}

\begin{itemize}
\item[$\square$] 9 月 18 日起登录系统注册，核对本人推免资格
\item[$\square$] 填报志愿前\textbf{通读目标单位发布的接收推免生章程和专业目录}
\item[$\square$] 记住 48 小时锁定规则，合理安排三个平行志愿的投递顺序
\item[$\square$] 收到复试通知后\textbf{在招生单位规定时间内确认}，逾期招生单位可取消
\item[$\square$] 收到待录取通知后\textbf{立即确认各校规定的确认时限}——各校差异极大，短的可能只有几十分钟
\item[$\square$] 确认前想清楚：接受待录取后\textbf{不能再填报其他志愿、不能再接受其他通知}
\end{itemize}

\subsection*{阶段四：录取之后（次年 6 \textasciitilde{} 9 月）}

\begin{itemize}
\item[$\square$] 完成本科培养方案要求的全部学分，\textbf{末学年不得有不及格}
\item[$\square$] 确保无违纪处分，毕业论文达到学校要求（部分学校要求优良或 80 分以上）
\item[$\square$] 按期取得\textbf{毕业证书与学位证书}
\item[$\square$] 按要求完成体检、政审与档案调取
\item[$\square$] 关注录取通知书发放与入学报到安排
\end{itemize}

\noindent 说明：以上清单为通用框架，\textbf{具体项目以你所在学校与录取单位当年的通知为准}。入学资格复审的全国性依据为《普通高等学校学生管理规定》（教育部令第 41 号）：报到时进行入学资格初步审查，材料与本人不符或违反招生考试规定者取消入学资格；入学后 3 个月内复查，弄虚作假、徇私舞弊者取消学籍。\nref{S0203}

\section{常见误区速查}

下表汇总本手册第 \ref{ch19} 章逐条核对的常见说法。\textbf{“正确方向”是判断方向，不是可以直接照搬的结论}；涉及具体门槛与名额的，必须回到你所在学校与目标学院的当年文件核实。

\begin{longtable}{p{6.4cm}p{9.6cm}}
\toprule
\textbf{流传的说法} & \textbf{实际方向} \\ \midrule
\endfirsthead
\toprule
\textbf{流传的说法} & \textbf{实际方向} \\ \midrule
\endhead
必须过六级才能保研 & 多数学校无校级硬性六级要求，门槛普遍下沉到学院／项目层面 \\
必须发论文才能保研 & 论文通常是加分项而非资格项，资格由综合测评排名决定 \\
有竞赛国奖就能保研 & 加分受比例上限与认定目录约束，不能替代学业排名 \\
排名第一想去哪去哪 & 排名只解决校内资格，外校录取取决于对方考核 \\
双非一定没机会保 985 & 本科出身影响初筛，进入考核环节后个人表现权重上升 \\
参加夏令营就有 offer & 多数营宽进严出，入营名单、优营名单、录取名单规模差数倍 \\
优营等于录取 & 优营效力分“直接录取／优先候补／仅荣誉”三档，差別巨大 \\
预推免只是走过场 & 不少院系的预推免录取占比已超过夏令营 \\
九推只能捡漏 & 大量正式名额在系统开放后释放，含优质项目 \\
系统只能填一个志愿 & 可同时填报三个平行志愿（专项计划另有规定者除外） \\
志愿填得越早越有优势 & 录取取决于招生单位审核发通知，与填报先后无关 \\
必须联系到导师才有戏 & 多数院系为委员会制，导师个人权重有限；不回复是常态 \\
确认待录取后绝对不能反悔 & 制度上有放弃路径，但代价逐年收紧，且会占用他人机会 \\
候补就是没戏 & 九推阶段候补转正并不罕见，主动跟进影响很大 \\
花钱买“保过”服务能稳上 & 不存在能操控录取结果的合法服务，此类宣传多为信息包装 \\
\bottomrule
\end{longtable}
